\chapter{Research Methods}

\section{Epistemological Foundations: A Priori and A Posteriori}

\begin{KeyConcept}
\begin{itemize}
\item \textbf{Epistemology} = the study of knowledge --- how you know what you know.
\item \textbf{A priori knowledge} --- knowledge \emph{before} facts, based on \textbf{intuition} (e.g. 'you cannot visit that haveli after 6 pm' $\rightarrow$ it must be haunted). It rests on \textbf{verifiability} (you keep verifying your own intuition rather than testing it).
\item \textbf{A posteriori knowledge} --- knowledge \emph{after} facts, based on \textbf{experiment/observation}, resting on the principle of \textbf{falsifiability}.
\end{itemize}
\end{KeyConcept}

\section{Quantitative and Qualitative Methods}

\begin{KeyConcept}
\begin{itemize}
\item \textbf{Quantitative methods} are used by \textbf{positivists} (who believe the social world can be theorised, with data, laws and causal relations); \textbf{qualitative methods} are used by \textbf{non-positivists}.
\item \textbf{Quantitative research methods}: \textbf{official statistics} (Durkheim's police-station suicide data), \textbf{survey}, \textbf{structured interview} and \textbf{questionnaire}.
\item \textbf{Qualitative research methods}: \textbf{ethnography} (with \textbf{participant observation}), the \textbf{comparative method}, \textbf{case study}, \textbf{life history} and \textbf{oral history} (a history transmitted orally through families, never recorded).
\item The method you choose depends on your \textbf{research question}: for the religion of the (uncontacted) \textbf{Sentinelese}, a survey/questionnaire/structured interview/oral history is not feasible (life at risk, no language); \textbf{ethnography or case study} (reading books/documentaries on them) is.
\item \textbf{Malinowski} ('Argonauts of the Western Pacific', the Trobriand Islanders) popularised participant observation; \textbf{Durkheim, by contrast, never did fieldwork} --- he collected data on the \textbf{Arunta} tribe from the ethnographers \textbf{Spencer and Gillen} (secondary data). \textbf{Margaret Mead} (Papua New Guinea) observed three tribes and found males feminine and females masculine in some --- showing 'males behave in masculine manner universally' is wrong.
\item For 'did the joint family break down due to industrialisation?' the \textbf{comparative method} (case study before vs after) is apt; for 'consumer behaviour of Delhi University students' the \textbf{survey} is apt.
\item \textbf{Quantitative methods are more reliable; qualitative methods are more valid.}
\item Even \textbf{religiosity and happiness can be quantified} --- the \textbf{Pew survey} measures how religious people are, just as there is a happiness index --- collecting both quantitative and qualitative data.
\end{itemize}
\end{KeyConcept}

\section{Triangulation and Primary/Secondary Data}

\begin{KeyConcept}
\begin{itemize}
\item \textbf{No contemporary research relies on one method} --- multiple methods are used (a mix of quantitative and qualitative), and this process is called \textbf{triangulation}.
\item \textbf{Primary data} --- data collected by you, first-hand (e.g. visiting a village and recording its population). \textbf{Secondary data} --- data collected by someone else which you interpret (e.g. crime/poverty data from newspapers); the same data is primary to the collector and secondary to the consumer.
\item \textbf{Techniques of data collection} include \textbf{participant observation} and \textbf{interviews} (the same methods are also techniques, since through them you collect data).
\end{itemize}
\end{KeyConcept}

\section{Variables and Hypothesis}

\begin{KeyConcept}
\begin{itemize}
\item A \textbf{variable} is any category \textbf{subject to change} (e.g. temperature, poverty, unemployment). The \textbf{independent variable} is the \textbf{cause}; the \textbf{dependent variable} is the \textbf{effect} ('poverty leads to unemployment' --- poverty = independent, unemployment = dependent).
\item \textbf{Durkheim} used \textbf{multivariate analysis} (Protestants/Catholics, men/women, aged/young) to show that the degree of integration/regulation determines suicide; \textbf{Marx} made the mode of production the dependent variable of the means and relations of production; \textbf{Weber} made capitalism the dependent variable of the Protestant ethic (but held a \textbf{multi-causal} world --- the Protestant ethic is only the most unique of several causes).
\item A \textbf{hypothesis} is a statement/assumption reflecting the \textbf{relation between variables}. \textbf{Simple} hypothesis = two variables (A $\rightarrow$ B); \textbf{complex} = many variables; \textbf{null hypothesis} = when no relation is found (so it is rejected and an \textbf{alternative hypothesis} is developed).
\item \textbf{Sources of a hypothesis}: literature review (Weber reviewing Benjamin Franklin), life history/biography, \textbf{dream experience} (Freud's 'Interpretation of Dreams'; Kekulé's benzene ring), theories (Marxism, feminism), documentaries/films/songs (Pink Floyd's 'We Don't Need No Education'), or simply a conversation.
\item \textbf{Importance of the hypothesis}: it directs the research, may lead to further hypotheses, and helps develop and test generalizations.
\end{itemize}
\end{KeyConcept}

\begin{QuickRevision}
\begin{itemize}
\item A priori (intuition, verifiability, before fact) vs a posteriori (experiment, falsifiability, after fact).
\item Quantitative (positivist): official statistics, survey, structured interview, questionnaire. Qualitative (non-positivist): ethnography/participant observation, comparative, case study, life history, oral history.
\item Method depends on research question (Sentinelese $\rightarrow$ ethnography; joint family $\rightarrow$ comparative; DU students $\rightarrow$ survey).
\item Triangulation = using multiple methods; primary vs secondary data.
\item Variable = category subject to change; independent (cause) vs dependent (effect); Durkheim multivariate, Marx (mode of production), Weber (multi-causal).
\item Hypothesis: simple/complex/null/alternative; sources (literature review, life history, dream, theory, songs, conversation); importance (directs, leads, tests).
\end{itemize}
\end{QuickRevision}

\section{Reliability and Validity}

\begin{KeyConcept}
\begin{itemize}
\item \textbf{Reliability} = the extent to which a result is \textbf{consistent when repeated}; \textbf{Validity} = the extent to which a method \textbf{measures what it intends to measure}.
\item \textbf{Positivists are concerned with reliability} (consistent findings); \textbf{non-positivists with validity} (in-depth, valid information). So \textbf{quantitative methods (structured interview, questionnaire) have high reliability}, while \textbf{qualitative methods yield more valid data}.
\item \textbf{Atkinson and J. D. Douglas} (ethnomethodologists) questioned the \textbf{validity} of Durkheim's suicide data --- what is 'suicide' to Durkheim might actually be \textbf{murder} (a validity issue, not reliability).
\item \textbf{The Aarushi Talwar case}: structured interviews --- two CBI teams asked the parents the same (scripted) questions; consistent answers = a \textbf{reliable} finding (only four people were in the house).
\item \textbf{Types of reliability} (names only): \textbf{test-retest} (one set at two points of time --- like rating Domino's on different platforms); \textbf{parallel forms} (one set, two points of time, two methods --- questionnaire vs the Likert scale); \textbf{inter-rater} (one set, one method, two assessors --- the two CBI teams, whose opposite findings were \emph{not} reliable, which is why the Supreme Court closed the case for lack of proof).
\end{itemize}
\end{KeyConcept}

\begin{QuickRevision}
\begin{itemize}
\item Reliability = consistent on repeat; Validity = measures what intended.
\item Positivists $\rightarrow$ reliability (quantitative: structured interview/questionnaire); non-positivists $\rightarrow$ validity (qualitative, in-depth).
\item Atkinson \& J. D. Douglas: Durkheim's 'suicide' may be murder $\rightarrow$ validity question.
\item Aarushi case $\rightarrow$ structured interview, reliable finding; types of reliability: test-retest, parallel forms (Likert scale), inter-rater (two CBI teams).
\end{itemize}
\end{QuickRevision}

\section{The Comparative Method}

\begin{KeyConcept}
\begin{itemize}
\item \textbf{Durkheim}: 'the comparative method is not just a method --- it is \textbf{sociology itself}' (no generalisation is possible without comparison).
\item Comparison can be done at \textbf{three levels}: (1) \textbf{two different societies at one point of time} (India 2020 vs China 2020); (2) \textbf{one society at two different points of time} (India 1947 vs today); (3) \textbf{different social groups within one society} at one point of time (Protestants vs Catholics).
\item Examples of the three levels: \textbf{Durkheim} (compared social groups --- Protestants/Catholics, married/unmarried, aged/young $\rightarrow$ degree of integration explains suicide; and compared pre-industrial vs industrial society $\rightarrow$ 'suicide increases with industrialisation'); \textbf{Weber} (compared the Protestant ethic with Hindu, Confucian and Judaic ethics $\rightarrow$ 'religious ideas are responsible for economic growth', \emph{turning Marx upside down}); \textbf{Marx} (compared ancient, feudal and capitalist modes of production $\rightarrow$ 'the history of all society is the history of class struggle', by inductive reasoning); \textbf{Talcott Parsons} (compared pattern variable A --- traditional --- with B --- modern; evolutionary universals); \textbf{Robert Merton} (compared social groups within 1940s America --- innovators, ritualists, rebels, retreatists $\rightarrow$ the strain theory); \textbf{Theda Skocpol} (compared the Chinese, American and French revolutions $\rightarrow$ for a revolution, two forces are needed: one decapitating the system from within, one attacking from outside); \textbf{Oliver Cox} (compared race with caste $\rightarrow$ 'race is coloured caste').
\item The comparative method \textbf{helps develop generalizations and even laws of human behaviour, and models of social change} (e.g. Marx's historical materialism).
\item \textbf{Alan Macfarlane}: instead of always comparing two (which breeds \textbf{relativism and essentialism} --- 'Scotland is better than Shimla' because the comparer is from Scotland; 'French secularism is the real secularism'), develop a \textbf{triadic model} (three or more) to avoid assuming one is superior. (Ashis Nandy: India had secularism before France --- Akbar's \textbf{Sulh-i-kul} and \textbf{Ibadat Khana}.)
\end{itemize}
\end{KeyConcept}

\begin{QuickRevision}
\begin{itemize}
\item Durkheim: 'comparative method is sociology itself'.
\item Three levels: two societies at one time; one society at two times; groups within one society.
\item Durkheim (integration/suicide), Weber (religion $\rightarrow$ economy; turned Marx upside down), Marx (class struggle; inductive), Parsons (pattern variables), Merton (strain), Skocpol (revolutions), Cox (race = coloured caste).
\item Macfarlane: triadic model to avoid relativism/essentialism (Ashis Nandy --- Akbar's Sulh-i-kul, Ibadat Khana).
\end{itemize}
\end{QuickRevision}

\section{Sampling}

\begin{KeyConcept}
\begin{itemize}
\item \textbf{Sampling} = the method of selecting a population. A \textbf{sample} should be \textbf{representative} and easy to collect.
\item \textbf{Random (probabilistic) sampling} --- everyone has an equal chance of selection: \textbf{simple random sampling} (slips in a basket); \textbf{stratified random sampling} (divide into strata --- gender, caste, religion --- and pick from each); \textbf{systematic random sampling} (a random starting point + a fixed sampling interval, e.g. every 10th candidate); \textbf{clustered random sampling}.
\item \textbf{Non-random (non-probabilistic) sampling} --- a part of the population is favoured: \textbf{convenience sampling} (whosoever is available, e.g. only Delhi candidates); \textbf{voluntary response sampling} (people volunteer themselves); \textbf{purposive sampling} (the researcher uses judgement to select); \textbf{snowball sampling} (contacts through contacts, e.g. reaching transgenders through one friend); \textbf{quota sampling}.
\item \textbf{Example}: a study of Uber drivers' monthly income --- fieldwork (primary data) at one petrol pump, collecting \textasciitilde{}60 drivers over 5 days, gave a \textbf{representative} (though not fully so --- mostly male, mostly Delhi) sample.
\end{itemize}
\end{KeyConcept}

\begin{QuickRevision}
\begin{itemize}
\item Sampling = selecting a population; sample should be representative and easy to collect.
\item Random: simple, stratified (strata), systematic (interval), clustered.
\item Non-random: convenience, voluntary response, purposive, snowball, quota.
\item Uber-driver petrol-pump example.
\end{itemize}
\end{QuickRevision}

